\documentclass[11pt,leqno]{article}
\headheight=13.6pt

% packages
\usepackage[alphabetic]{amsrefs}
\usepackage{physics}
% margin spacing
\usepackage[top=1in, bottom=1in, left=0.5in, right=0.5in]{geometry}
\usepackage{amsfonts, amsmath, amssymb, amsthm}
\usepackage{extarrows}
\usepackage[none]{hyphenat}
\usepackage{fancyhdr}
\usepackage{graphicx}
\graphicspath{{./images/}}
\usepackage{float}
\usepackage{color}
\newcommand{\sai}[1]{\textcolor{red}{#1}}
\usepackage{enumitem}
% \usepackage{mathrsfs}
\usepackage{hyperref}
\usepackage[noabbrev, capitalise]{cleveref}
\crefformat{equation}{equation~#2#1#3}
\crefformat{lemma}{\textrm{Lemma}~#2#1#3}
\usepackage{quiver}

% theorems
\theoremstyle{plain}
\newtheorem{lem}{Lemma}
\newtheorem{lemma}[lem]{Lemma}
\newtheorem{thm}[lem]{Theorem}
\newtheorem{theorem}[lem]{Theorem}
\newtheorem{prop}[lem]{Proposition}
\newtheorem{proposition}[lem]{Proposition}
\newtheorem{cor}[lem]{Corollary}
\newtheorem{corollary}[lem]{Corollary}
\newtheorem{conj}[lem]{Conjecture}
\newtheorem{fact}[lem]{Fact}
\newtheorem{form}[lem]{Formula}

\theoremstyle{definition}
\newtheorem{defn}[lem]{Definition}
\newtheorem{definition/}[lem]{Definition}
\newenvironment{definition}
  {\renewcommand{\qedsymbol}{\textdagger}%
   \pushQED{\qed}\begin{definition/}}
  {\popQED\end{definition/}}
\newtheorem{example}[lem]{Example}
\newtheorem{remark}[lem]{Remark}
\newtheorem{exercise}[lem]{Exercise}
\newtheorem{notation}[lem]{Notation}

\numberwithin{equation}{section}
\numberwithin{lem}{section}

% header/footer formatting
\pagestyle{fancy}
\fancyhead{}
\fancyfoot{}
\fancyhead[L]{Invariant theory and geometry}
\fancyhead[C]{}
\fancyhead[R]{Sai Sivakumar}
\fancyfoot[R]{\thepage}
\renewcommand{\headrulewidth}{1pt}

% paragraph indentation/spacing
\setlength{\parindent}{0cm}
\setlength{\parskip}{10pt}
\renewcommand{\baselinestretch}{1.25}

% mathscript S
\usepackage[mathscr]{euscript}

% math operators
\DeclareMathOperator{\Stab}{Stab}
\DeclareMathOperator{\Ind}{Ind}
\DeclareMathOperator{\Fun}{Fun}
\DeclareMathOperator{\id}{id}
\DeclareMathOperator{\End}{End}
\DeclareMathOperator{\GL}{GL}
\DeclareMathOperator{\SL}{SL}
\DeclareMathOperator{\SO}{SL}
\DeclareMathOperator{\Sp}{SL}
\DeclareMathOperator{\Lie}{Lie}
\DeclareMathOperator{\diag}{diag}

% bracket notation for inner product
\usepackage{mathtools}

\DeclarePairedDelimiterX{\abr}[1]{\langle}{\rangle}{#1}

% smileys frownies
\usepackage{wasysym}
\newcommand{\smallhappy}{\raisebox{-.14em}{\smiley}}
\newcommand{\happy}{\raisebox{-.24em}{\resizebox{1.2em}{!}{\smiley}}}
\newcommand{\smallsad}{\raisebox{-.14em}{\frownie}}
\newcommand{\sad}{\raisebox{-.24em}{\resizebox{1.2em}{!}{\frownie}}}
\DeclareMathOperator{\mathhappy}{\!\happy\!}
\DeclareMathOperator{\smallmathhappy}{\!\smallhappy\!}
\DeclareMathOperator{\mathsad}{\!\sad\!}
\DeclareMathOperator{\smallmathsad}{\!\smallsad\!}

% set page count index to begin from 1
\setcounter{page}{1}

% array column and row separation
\arraycolsep = 2pt
\renewcommand{\arraystretch}{.6}

\begin{document}
These are notes which abridge and expand on three talks I gave in the course M 392C Geometry Literature which Dr. Bernd Siebert organized in Fall 2026. He encouraged students in this course to use artificial intelligence to aid in preparation for these talks. I asked GPT-5.6 Sol and GPT-6 Astra various questions and got some help that way, but what appears in these notes is my own writing and appears originally from various texts.

\begin{abstract}
  \sai{to write}
\end{abstract}

\section{The premise}
Throughout, we work over the field of complex numbers. Consider $\mathbb A^n$ and its algebra of regular functions, which we identify with $R = \mathbb C[x_1,\dots,x_n]$. The group $\GL_n(\mathbb C)$ acts on $\mathbb A^n$ on the left by matrix multiplication, from which we obtain the left regular representation of $\GL_n(\mathbb C)$ on $R$ by
\[(A\cdot f)(x) = f(A^{-1}x)\quad \text{for}\quad x\in\mathbb A^n.\]
In the above we may replace $\GL_n(\mathbb C)$ by subgroups $G\subseteq \GL_n(\mathbb C)$.

Are there functions $f\in R$ for which $A\cdot f = f$ for all $A\in G$? Such functions are called $G$-invariants (suppressing the $G$ when it's understood). If they exist, how do we calculate them? The $\mathbb C$-algebra of $G$-invariants is denoted by $R^G$. Is $R^G$ finitely generated? If so, $R^G\cong \mathbb C[y_1,\dots,y_s]/J$ for some $s$ and $J$, and one might wonder how to calculate $J$.

These questions, among others, would lead to the development of classical invariant theory. In what follows, we will refine and answer some of these questions, mention some historical developments, and naturally migrate towards some ideas in geometric invariant theory and symplectic geometry.

Some examples: \begin{enumerate}
  \item When $G = S_n$, the symmetric polynomials are the $S_n$-invariants.
  \item For $U = \{\big(\!\begin{smallmatrix}
        1 & \ast \\ 0 & 1
  \end{smallmatrix}\!\big)\}$, $\mathbb C[x,y]^U = \mathbb C[y]$.
  \item When $G = \GL_n(\mathbb C)$, $\mathbb C[x_1,\dots,x_n]^{\GL_n(\mathbb C)} = \mathbb C$.
\end{enumerate}

A more refined situation we should consider is an algebraic action of an algebraic subgroup $G\subseteq \GL_n(\mathbb C)$ on $\mathbb A^m$ where $n$ need not equal $m$, or even on an affine variety $V\subseteq \mathbb A^m$. Like before, it makes sense to pass to the left regular representation on the algebra of regular functions $R = \mathbb C[V]$ and again ask about invariants.

Two examples: \begin{enumerate}
  \item The action of $\mathbb C^\times = \GL_1(\mathbb C)$ on $\mathbb C^2$ by $t\cdot \big(\!\begin{smallmatrix}
        x \\ y 
  \end{smallmatrix}\!\big) = \big(\!\begin{smallmatrix}
        t & 0 \\ 0 & t^{-1}
  \end{smallmatrix}\!\big)\big(\!\begin{smallmatrix}
        x \\ y 
  \end{smallmatrix}\!\big)$ yields the ring of invariants $\mathbb C[x,y]^{\mathbb C^\times} = \mathbb C[xy]$.
  \item The action of $\SL_n(\mathbb C)$ on the space of $n\times n$ complex matrices by left multiplication yields the ring of invariants $\mathbb C[\det]$, where $\det$ is the determinant function in $R = \mathbb C[x_{11},\dots,x_{nn}]$. Indeed, for $f\in R^{\SL_2(\mathbb C)}$, define a function $p\in \mathbb C[t]$ by
  \[p(t) = f(\diag(t,1,\dots,1))\]
  and for any $A\in \GL_n(\mathbb C)$, $f(A) = p(\det(A))$ since $A = g\diag(\det(A),1,\dots,1)$ for some $g\in \SL_2(\mathbb C)$ and $f$ is an $\SL_2(\mathbb C)$-invariant. Then by Zariski density of $\GL_n(\mathbb C)$ in the space of $n\times n$ complex matrices, it follows that $f(A) = p(\det(A))$ for any complex $n\times n$ matrix $A$.
\end{enumerate}

From now on, we will consider only reductive complex groups, which are a nice collection of algebraic subgroups of $\GL_n(\mathbb C)$ for some $n$. Examples of complex reductive groups are ``groups with a name''; for example, $\GL_n(\mathbb C)$, $\SL_n(\mathbb C)$, $(\mathbb C^\times)^n$, $\SO_n(\mathbb C)$, $\Sp_{2n}(\mathbb C)$, and finite groups, to name a few.

\section{Finiteness of invariants}
\sai{statement of theorem, unitary trick, Hilbert basis theorem (modern and Hilbert's formulations but no proof), gesture at 1893 algorithms but these should differ from Gr\"obner basis calculations}

\section{Syzygies of invariants}
\sai{modern and Hilbert 1893 statements (not proofs) of Syzygy theorem and explain how one could build these free resolutions by hand, what they mean etc. Can gesture at why homological methods in commutative algebra are important cf Eisenbud sequel to comm. alg. w/ view towards alg. geo. and Erman, Berkesch and friends' program (see JMM 2024 Erman talk). Could give small example of Hilbert-Burch resolution for two variables}

\section{The variety corresponding to the ring of invariants}
\sai{in what sense is $V(J)$ a quotient of $\mathbb A^m$? Wrestle with orbit-closure details and state relevant theorems which make sense of or answer this question. Give some examples and explain a little about the subset of stable points and how restriction of the projection map actually is a quotient. But the quotient on the whole affine space results in what is called the geometric invariant theory quotient (GIT quotient)}

\section{The Nullcone, projective varieties}
\sai{a complete mystery, to be honest}

\section{Symplectic geometry and the Kempf-Ness theorem}
\sai{talk 3, to be determined}

\end{document}